\section*{Capítulo \ref{ch:b2:phylogenetic-trees} --- Árboles filogenéticos y reloj molecular}

\begin{solution}{exo:b2:phylogenetic-trees:1}
\omterm{def:b2:phylogenetic-trees:tree}{Homología}: un carácter heredado de un antepasado común --- el húmero, el
radio y el cúbito del ala de un murciélago y de la aleta de una ballena.
\omterm{def:b2:phylogenetic-trees:tree}{Analogía}: un carácter parecido adquirido de forma independiente --- las
alas de los murciélagos y de las aves como superficies de vuelo (los
huesos son homólogos como miembros anteriores; el ala como tal no lo es),
o el ojo en cámara del pulpo y el de los vertebrados. \omterm{def:b2:phylogenetic-trees:tree}{Sinapomorfía}: una
\omterm{def:b2:phylogenetic-trees:tree}{homología} derivada compartida por un \omterm{def:b2:phylogenetic-trees:tree}{clado} y que lo define --- las plumas
para las aves, el huevo amniótico para los amniotas.
\end{solution}

\begin{solution}{exo:b2:phylogenetic-trees:2}
Aves: monofilético. Reptiles sin las aves: parafilético. Peces sin los
tetrápodos: parafilético (los peces pulmonados están más próximos a
nosotros que a la trucha). Vertebrados voladores (murciélagos, aves,
pterosaurios): polifilético. Mamíferos: monofilético. \omterm{def:b2:unicellular-diversity:prokaryote}{Procariotas}:
parafilético (bacterias y arqueas, dejando fuera a los eucariotas --- y las
arqueas están más próximas a los eucariotas que a las bacterias).
\end{solution}

\begin{solution}{exo:b2:phylogenetic-trees:3}
Grupo hermano de las aves: los cocodrilos. Grupo hermano de los
mamíferos: todos los demás amniotas (tortugas, lagartos, cocodrilos y aves
juntos). \omterm{def:b2:phylogenetic-trees:tree}{Clado} más pequeño que contiene a los lagartos y las aves:
(lagartos,(cocodrilos, aves)), los diápsidos sin las tortugas en este
árbol. Intercambiar los cocodrilos y las aves en su nodo no cambia nada:
un nodo puede girarse libremente.
\end{solution}

\begin{solution}{exo:b2:phylogenetic-trees:4}
No enraizados: $3$; $3\times 5\times 7 = 105$; $3\times 5\times
7\times 9\times 11\times 13\times 15 = \num{2027025}$. Árboles enraizados
para 4 taxones: cada uno de los 3 árboles no enraizados tiene 5 ramas, lo
que da 15.
\end{solution}

\begin{solution}{exo:b2:phylogenetic-trees:5}
$((A,B),(C,D))$: sitios 1, 2 y 3, un paso cada uno; sitio 4, dos pasos
($\mathrm{A}$ y $\mathrm{G}$ a ambos lados); sitio 5, uno: 6 pasos.
$((A,C),(B,D))$: $2 + 2 + 2 + 1 + 1 = 8$. $((A,D),(B,C))$: $2 + 2 + 2 +
2 + 1 = 9$. El primer árbol es el más parsimonioso; en él, el sitio 4 es
homoplásico (el cambio $\mathrm{A}\to\mathrm{G}$, o su inverso, ocurre dos
veces).
\end{solution}

\begin{solution}{exo:b2:phylogenetic-trees:6}
Se unen $A$ y $B$ a la altura 3. Entonces $d(AB,C) = 14$, $d(AB,D) = 16$,
$d(C,D) = 10$: se unen $C$ y $D$ a la altura 5. Entonces $d(AB,CD) = (14 + 16
+ 14 + 16)/4 = 15$: raíz a la altura $7.5$. Árbol $((A,B),(C,D))$.
\end{solution}

\begin{solution}{exo:b2:phylogenetic-trees:7}
$d = -0.75\ln(1 - 4p/3)$: $0.052$, $0.38$, $1.21$. Con $p = 0.6$, el
argumento del logaritmo es $0.2$ y la pendiente $\mathrm{d}d/\mathrm{d}p
= 1/(1 - 4p/3) = 5$: un error de muestreo de $0.02$ en $p$ se convierte en
$0.1$ en $d$, y la hipótesis de tasas iguales en todos los sitios, falsa en
cualquier gen real, añade un sesgo del mismo orden. El gen está saturado.
\end{solution}

\begin{solution}{exo:b2:phylogenetic-trees:8}
$t = d/2k = 0.02/(2\times 10^{-9}) = 10$ millones de años. A 500
millones de años, $d = 2\times 10^{-9}\times 5\times 10^{8} = 1.0$ y
$p = 0.75(1 - e^{-4/3}) = 0.55$: tres cuartas partes del camino hacia la
saturación, donde la corrección es muy empinada y la fecha poco fiable
--- para esas separaciones hace falta un gen más lento.
\end{solution}

\begin{solution}{exo:b2:phylogenetic-trees:9}
\qty{52}{\%}: el \omterm{def:b2:phylogenetic-trees:tree}{clado} se recupera en apenas la mitad de los remuestreos;
los datos casi no dicen nada sobre él, y las alternativas juntas son
igual de probables. \qty{99}{\%}: la señal es coherente en todos los
sitios; el \omterm{def:b2:phylogenetic-trees:tree}{clado} está fuertemente respaldado (dados el modelo y el
alineamiento --- el bootstrap no detecta errores sistemáticos como la
\omterm{prop:b2:phylogenetic-trees:pitfalls}{atracción de ramas largas}). Para el primero, añádanse secuencias (más
genes) y taxones que rompan las ramas, y compruébese con otro método.
\end{solution}

\begin{solution}{exo:b2:phylogenetic-trees:10}
Dos linajes que han cambiado cada uno en una gran fracción de sitios
tienen, en cualquier sitio donde ambos han cambiado, una probabilidad de
uno entre cuatro de llegar a la misma base (con cuatro bases a tasas
iguales). Si cada uno ha cambiado en el \qty{40}{\%} de los sitios, ambos
han cambiado en el \qty{16}{\%} y coinciden por azar en el \qty{4}{\%} de
todos los sitios --- lo que puede superar la fracción de sitios que portan
las \omterm{def:b2:phylogenetic-trees:tree}{sinapomorfías} verdaderas que unen a cada uno con sus parientes reales,
de evolución lenta. La parsimonia cuenta coincidencias sin preguntarse por
qué se producen y une a los dos. Añadir taxones que se ramifiquen a lo
largo de cada rama larga la divide en segmentos más cortos, de modo que
las coincidencias fortuitas se reparten entre varias ramas más cortas y las
\omterm{def:b2:phylogenetic-trees:tree}{sinapomorfías} verdaderas de cada segmento se hacen visibles; un modelo de
verosimilitud que espera muchas coincidencias fortuitas en las ramas
largas las corrige directamente.
\end{solution}

\begin{solution}{exo:b2:phylogenetic-trees:11}
$d = -0.75\ln(1 - 0.12) = 0.096$; $t = d/2k = 0.096/(4\times 10^{-8})
= 2.4$ millones de años. Sustituciones: $0.096\times \num{16000} =
1500$, precisión relativa $1/\sqrt{1500} = \qty{2.6}{\%}$: un error
estadístico de $\pm 0.06$ millones de años. Pero la tasa está calibrada
sobre otras separaciones y puede ser errónea en un factor dos para este
linaje (variación de tasas entre linajes, y diferencia entre la tasa
genealógica medida a lo largo de generaciones y la tasa medida a lo largo
de millones de años); y la divergencia del gen es anterior a la de las
especies, puesto que la población ancestral era polimórfica --- de modo que
la verdadera separación es más reciente que la del gen, en un tiempo que
depende del tamaño de la población ancestral. (La fecha aceptada, a partir
de muchos genes nucleares y de fósiles, es de 6 a 7 millones de años: el
reloj mitocondrial va aquí demasiado deprisa.)
\end{solution}

\begin{solution}{exo:b2:phylogenetic-trees:12}
\omterm{prop:b2:phylogenetic-trees:pitfalls}{Separación incompleta de linajes}: cuando dos especiaciones están próximas
en el tiempo, un polimorfismo ancestral puede repartirse de modo que el
árbol de un gen agrupe las especies de forma distinta al árbol de
especies; transferencia horizontal: un gen recibido de otro linaje lleva
la historia de ese linaje; \omterm{prop:b2:genome-diversification:duplication}{duplicación génica}: comparar un parálogo de una
especie con el ortólogo de otra da la fecha de la duplicación, no la de la
especiación. Un árbol de especies se obtiene usando muchos genes de todo
el genoma y tomando el árbol que respalda la mayoría (o un modelo que
espera una fracción conocida de genes discordantes), eligiendo con cuidado
los ortólogos y prefiriendo, en los \omterm{def:b2:unicellular-diversity:prokaryote}{procariotas}, los genes ribosómicos e
informacionales, que rara vez se transfieren.
\end{solution}

\begin{solution}{pb:b2:phylogenetic-trees:1}
\textbf{1.} Sitios 1, 2 y 4 (estados compartidos solo por $W$ y $H$,
derivados respecto del camello): \omterm{def:b2:phylogenetic-trees:tree}{sinapomorfías} de $\{W,H\}$. Sitios 3
y 6: \omterm{def:b2:phylogenetic-trees:tree}{sinapomorfías} de $\{W,H,C\}$. El sitio 5 es una autapomorfía de
$W$, un cambio en un solo linaje que no agrupa nada; ningún sitio es
homoplásico en el árbol molecular.
\textbf{2.} Un cambio por sitio: 6 pasos.
\textbf{3.} Los sitios 1, 2 y 4 necesitan ahora dos cambios cada uno, y los
sitios 3, 5 y 6, uno: 9 pasos.
\textbf{4.} El árbol ballena--hipopótamo, por tres pasos. Tres sitios de
cada mil son un margen pequeño que unas pocas homoplasias podrían
invertir; reforzarlo supone más genes, más taxones (otros rumiantes,
pecaríes, ciervos) y un bootstrap.
\textbf{5.} El \omterm{def:b2:phylogenetic-trees:tree}{clado} reaparece en cinco remuestreos de cada seis: un
respaldo moderado, muy por encima del azar pero por debajo del
\qty{95}{\%} que se considera convencionalmente fuerte.
\textbf{6.} El \omterm{meth:b2:phylogenetic-trees:parsimony}{grupo externo} debe estar fuera del grupo, pero lo bastante
cerca para que su secuencia siga siendo alineable y no esté saturada; el
camello es un ungulado de dedos pares situado fuera del grupo
vaca--cerdo--hipopótamo--ballena. Un pez diferiría a niveles casi saturados,
sus estados no informarían sobre qué estado es ancestral dentro de los
mamíferos, y su rama larga atraería al linaje del grupo interno de
evolución más rápida.
\textbf{7.} $d = 0.0625$, $0.107$, $0.131$, $0.155$.
\textbf{8.} Se unen $W$ y $H$ a la altura $0.031$. Entonces $d(WH,C) =
0.107$, $d(WH,P) = 0.131$, $d(C,P) = 0.131$, y todo a $M$,
$0.155$: se unen $WH$ y $C$ a $0.054$. Entonces $d(WHC,P) = 0.131$: unión
a $0.065$. Por último se une $M$ a $0.078$. Árbol $(M,(P,(C,(W,H))))$.
\textbf{9.} Sí: el mismo árbol, el mismo anidamiento.
\textbf{10.} Una ballena que hubiera cambiado el doble de deprisa estaría
más lejos de todo, incluido el hipopótamo; el UPGMA, que coloca cada punta
a la altura cero y une el par más próximo, uniría primero el hipopótamo
con la vaca y dejaría fuera a la ballena --- exactamente el árbol de los
anatomistas, por una razón equivocada.
\textbf{11.} $0.0625\times 1000 = 62.5$ sustituciones; desviación típica
de Poisson, $\sqrt{62.5} = 7.9$.
\textbf{12.} $7.9/62.5 = \qty{13}{\%}$.
\textbf{13.} $k = d/2t = 0.107/(1.2\times 10^{8}) = 8.9\times 10^{-10}$
por sitio por año.
\textbf{14.} $t = 0.0625/(2\times 8.9\times 10^{-10}) = 35$ millones
de años.
\textbf{15.} Cerdo: $0.131/(1.79\times 10^{-9}) = 73$ millones de años;
camello: $0.155/(1.79\times 10^{-9}) = 87$ millones de años.
\textbf{16.} $35 \pm 4.5$ millones de años (\qty{13}{\%}).
\textbf{17.} La tasa varía como $1/t_{\text{cal}}$ y la fecha como
$t_{\text{cal}}$: $\pm 5/60 = \qty{8}{\%}$, es decir, $\pm 3$ millones de
años --- combinado con el error de Poisson, unos $\pm 5$ millones de años.
\textbf{18.} $d = -0.75\ln(1 - 0.6) = 0.69$: el gen ha cambiado en la
práctica en dos tercios de sus sitios, la corrección ha multiplicado
$p$ por $1.5$ y su pendiente es $2.5$; está cerca de la saturación para
esta separación y su fecha vale poco.
\textbf{19.} La ausencia de un carácter no es un estado derivado
compartido: las ballenas perdieron por completo la pata trasera y su
astrágalo, de modo que no puede decirse nada sobre qué astrágalo habrían
tenido. Una \omterm{def:b2:phylogenetic-trees:tree}{sinapomorfía} es una prueba de un \omterm{def:b2:phylogenetic-trees:tree}{clado}; su pérdida en un linaje
no es una prueba en contra de la pertenencia.
\textbf{20.} Las ballenas fósiles tienen el astrágalo de doble polea: las
ballenas descienden de un antepasado que lo tenía, lo que las sitúa dentro
de los ungulados de dedos pares. Las rocas confirman la predicción
molecular.
\textbf{21.} Parafilético: un antepasado con todos sus descendientes
excepto la línea de las ballenas.
\textbf{22.} La \omterm{prop:b2:phylogenetic-trees:pitfalls}{separación incompleta de linajes} de un polimorfismo
ancestral a través de dos especiaciones rápidas (la vaca, el hipopótamo y
la ballena se separaron en unos pocos millones de años), o la comparación
de parálogos. Se decide secuenciando muchos genes y tomando el árbol que
respaldan la mayoría de ellos y su concatenación; y comprobando las
duplicaciones en cada familia.
\textbf{23.} $d = -0.75\ln(1 - 0.333) = 0.30$; $k = d/2t =
0.30/(7\times 10^{7}) = 4.3\times 10^{-9}$ por sitio por año, cinco
veces la del gen nuclear.
\textbf{24.} El ADN mitocondrial lo replica una polimerasa con una
corrección de errores más débil, está expuesto a los radicales de oxígeno
de la cadena respiratoria y no lo repara la maquinaria nuclear; además
carece de recombinación, de modo que los daños se acumulan.
\textbf{25.} Árbol $(M,(P,(C,(W,H))))$: las ballenas son el grupo hermano
de los hipopótamos; $k = 8.9\times 10^{-10}$ por sitio por año; divergencia
entre ballena e hipopótamo: $35 \pm 5$ millones de años.
\end{solution}
